\begin{abstract}
We characterize optimal estimation and honest inference for the population average treatment effect in a balanced randomized experiment with incomplete primary outcomes. Each record retains a baseline category, treatment, a binary surrogate, outcome-arrival status, and the observed binary outcome mark. Outcome arrival is missing at random conditional on baseline, treatment, and surrogate, with a known cellwise arrival floor \(q\in[1/2,1]\). For sample size \(n\ge1\) and at most \(d\ge1\) baseline categories, the minimax squared-error risk is, up to universal constants,
\[
\frac1n+(1-q)^2\min\left\{1,\left[\frac{d}{n\log(e+n)}\right]^2\right\}.
\]
For each fixed noncoverage level \(\alpha\in(0,1/2)\), the minimax worst-law expected length of honest connected confidence intervals has the square-root order of this benchmark. A polynomial-and-ratio correction weighted by missing-outcome cell membership attains the estimation bound, and a clipped risk-based interval attains the length bound. Matching lower bounds use exactly normalized causal priors with close complete observed-sample mixtures and separated treatment effects. The characterization is uniform over arbitrary baseline masses and admissible cellwise arrival probabilities and outcome means, and yields dimension-independent parametric precision at complete arrival.
\end{abstract}

\section{Introduction}

This paper characterizes how outcome arrival and baseline dimension jointly determine optimal precision for a population average treatment effect in a balanced randomized experiment. The experiment records baseline category, treatment assignment, a binary surrogate measurement, and outcome-arrival status for every unit, together with a binary primary outcome when it arrives. Arrival is missing at random conditional on baseline category, treatment, and surrogate. Under a known cellwise arrival floor between one half and one, we establish matching finite-sample bounds for minimax squared-error risk and for the worst-law expected length of connected confidence intervals with uniform coverage.

The model permits arbitrary baseline masses, including many sparsely occupied categories, and freely varying cellwise arrival probabilities and outcome means subject to its causal and positivity restrictions. Treatment is an independent fair coin, and realized surrogate measurements and outcomes agree with their assigned potential measurements. The population target is the mean difference between the two potential primary outcomes. These conditions define the randomized-MAR law class in \cref{obj:def:law-class}; the corresponding observed-law identification formula is established in \cref{obj:lem:identification-infrastructure}. The resulting uniform estimation problem concerns recovery of an aggregate treatment effect across cells whose individual outcome means may be sparsely observed.

Let \(n\ge1\) denote the sample size, \(d\ge1\) the baseline-support bound, and \(q\in[1/2,1]\) the known outcome-arrival floor. Our principal result, \cref{obj:thm:point-frontier}, establishes that minimax squared-error risk over this class is comparable to
\[
\frac1n+(1-q)^2
\min\left\{1,\left[\frac{d}{n\log(e+n)}\right]^2\right\},
\]
with universal comparison constants uniform over all such \(n,d,q\). For each fixed noncoverage level \(\alpha\in(0,1/2)\), \cref{obj:thm:interval-frontier} establishes that the minimax worst-law expected length among honest connected confidence intervals is comparable to the square root of this expression, with constants depending only on \(\alpha\). Honesty here requires coverage at least \(1-\alpha\) under every law in the specified class at each sample size.

The rate separates sampling uncertainty from the cost of recovering outcome means across many baseline categories. The additional root mean squared error is proportional to the arrival deficit \(1-q\), multiplied by the smaller of one and \(d/[n\log(e+n)]\). Thus improving arrival permits larger baseline supports at sampling-order precision. In particular, an arrival deficit at most \(n^{-1/2}\) yields root mean squared error and optimal honest expected length of order \(n^{-1/2}\), uniformly over the support bound. At complete arrival, \cref{obj:prop:complete-arrival-reduction} gives dimension-independent parametric rates and identifies the estimator as a projected outcome contrast using the known treatment-allocation weights.

The constructive argument starts from a simple decomposition: arrived outcomes estimate the observed part of each cell contribution, while records with missing outcomes still reveal their cell membership and therefore determine the mass requiring correction. The estimator in \cref{obj:def:mm-estimator} combines these two pieces. Separate auxiliary streams select between polynomial correction in lightly populated cells and an empirical ratio in heavily populated cells. Conditional averaging over the auxiliary randomization yields a procedure determined by the observed sample, and \cref{obj:lem:upper-risk} supplies its uniform squared-risk bound. A clipped interval based on this bound satisfies finite-sample uniform coverage and attains the expected-length order, as established in \cref{obj:lem:honest-interval-construction}. The exhibited polynomial tuning begins when \(\log(e+n)\ge128\).

The converse has two branches. When the squared arrival-weighted dimension scale is at least \(n^{-1}\), oppositely oriented pairs of baseline labels and a fixed filler label give exactly normalized priors whose complete observed-sample mixtures are close while their treatment effects concentrate on separated sides of a deterministic center; this is the comparison in \cref{obj:thm:normalized-converse}. When the squared arrival-weighted dimension scale is below \(n^{-1}\), the parametric floors supply the required lower bounds. Testing reductions combine these existential prior comparisons with parametric sampling floors to obtain both minimax lower bounds while preserving the factor \(1-q\).

The closest identification and efficiency benchmark is \citet[Lemma~1.1 under Assumptions~1--3, Theorem~2.1, and Sections~3--4]{KallusMao2025}, which develops surrogate-assisted ATE identification, efficiency, and estimation under its nuisance-convergence conditions. Our analysis characterizes uniform precision over unrestricted finite-cell nuisance functions under balanced randomization. The dimension-dependent complete-record analysis of \citet[Theorems~1--2 and Appendix~C.5]{ZengBalakrishnanHanKennedy2024} provides a related benchmark with unknown categorywise treatment propensities. Polynomial estimation for discrete treatment effects and an independent annotation experiment are developed in \citet{CausalSmith2026,CausalSmith2026Annotation}. The polynomial and moment-matching methods also build on \citet[Sections~III--IV]{JiaoVenkatHanWeissman2015} and \citet[Proposition~3 and Appendix~B.2]{WuYang2016}. Against these precedents, the present contribution is the joint arrival-and-dimension characterization for estimation and honest expected length in the specified randomized-MAR experiment.

The exposition proceeds through related work, the statistical setup and criteria, the minimax estimation and honest-inference results, and their statistical implications, followed by limitations and future work. The appendices develop identification and estimation bounds, the normalized-prior converse, and the testing reductions completing the main comparisons, together with the verification note and proofs of the main results.
